\chapter{Úvod}

\indent V súčasnosti zohrávajú úlohy obchodného cestujúceho (TSP) významnú úlohu v oblasti optimalizácie. S podobnými problémami sa stretávame v mnohých praktických situáciách a existuje veľké množstvo prístupov, ktoré umožňujú riešiť TSP aj jeho rôzne varianty. V praxi sa často používajú najmä heuristické a metaheuristické metódy, ktoré spravidla poskytujú suboptimálne riešenia.

V tejto práci sme sa na danú problematiku zamerali z pohľadu moderných metód umelej inteligencie. Motiváciou bolo overiť, či je možné využiť neurónové siete aj pri riešení TSP tak, aby boli výsledky použiteľné nielen z teoretického hľadiska, ale aj v praktických úlohách. Vychádzali sme z predpokladu, že aj napriek prirodzeným obmedzeniam môžu tieto metódy priniesť použiteľné riešenia a nové poznatky.

V rámci práce sme realizovali kompletný postup od generovania údajov, cez prípravu datasetov a trénovanie modelov, až po ich vyhodnotenie na referenčných benchmarkových dátach TSPLIB. Týmto postupom sme overili, že zvolený prístup dokázal riešiť úlohy TSP s merateľnou presnosťou a umožnil systematicky porovnať silné a slabé stránky jednotlivých modelových variantov.

Hlavnou motiváciou práce bolo tiež to, že umelá inteligencia sa dnes intenzívne zavádza do rôznych oblastí praxe. Preto sme považovali za dôležité overiť jej potenciál aj v optimalizačných úlohách tohto typu. Predložená práca tak môže slúžiť ako základ pre ďalší výskum a zároveň ako praktický rámec pre aplikáciu neurónových sietí pri riešení problémov TSP.

Cieľom práce bolo analyzovať možnosti riešenia optimalizačných úloh rôznymi prístupmi umelej inteligencie so zameraním na symetrický problém obchodného cestujúceho. Predmetom skúmania boli modely neurónových sietí určené na predikciu hrán v grafe, trénované na synteticky generovaných dátach a hodnotené na referenčných inštanciách TSPLIB. V práci sme sa sústredili na porovnanie kvality riešení pomocou optimalizačnej odchýlky voči referenčným trasám.

Prácu sme rozčlenili do viacerých logicky previazaných častí. V druhej kapitole sme predstavili teoretické východiská problematiky TSP, základné súvislosti umelej inteligencie a analýzu súčasného stavu riešenej problematiky. V tretej kapitole sme formulovali cieľ práce a podrobne opísali použitú metodiku, postup realizácie práce a nástroje využité pri experimentoch. Vo štvrtej kapitole sme uviedli dosiahnuté výsledky, ich porovnanie a diskusiu. V piatej kapitole sme zhrnuli hlavné závery práce a navrhli odporúčania pre ďalší výskum.
